%
% agent-identity-spec.tex
%
% Agent Identity Specification for Claude Code Sub-Agents
% How ethos generates, delivers, and maintains agent identity context
% across the Claude Code sub-agent lifecycle.
%
% Compile: pdflatex agent-identity-spec.tex
%

\documentclass[a4paper,10pt]{article}
\usepackage[margin=1in]{geometry}
\usepackage{booktabs}
\usepackage{listings}
\usepackage{xcolor}
\usepackage{enumitem}
\usepackage{longtable}
\usepackage{tabularx}
\usepackage[colorlinks=true,linkcolor=blue,citecolor=blue,urlcolor=blue]{hyperref}

\lstset{
  basicstyle=\ttfamily\small,
  keywordstyle=\color{blue!80!black},
  commentstyle=\color{green!50!black},
  stringstyle=\color{red!60!black},
  breaklines=true,
  frame=single,
  framerule=0.3pt,
  rulecolor=\color{gray!50},
  backgroundcolor=\color{gray!5},
  columns=fullflexible,
}

\newcommand{\fpath}[1]{\texttt{\hyphenchar\font=`\/ #1}}
\newcommand{\field}[1]{\texttt{#1}}
\newcommand{\cmd}[1]{\texttt{#1}}
\newcommand{\coderef}[1]{\textsf{\small [#1]}}

\tolerance=400
\emergencystretch=2em

\begin{document}

\title{Agent Identity Specification\\
\large How Ethos Generates and Delivers Identity\\
to Claude Code Sub-Agents}
\author{Punt Labs}
\date{April 2026 --- v3.0}
\hypersetup{
  pdftitle={Agent Identity Specification},
  pdfauthor={Punt Labs},
  pdfsubject={Agent Identity},
  pdfcreator={pdflatex}
}
\maketitle

\tableofcontents
\newpage

% ===================================================================
\section{Problem Statement}
% ===================================================================

Claude Code sub-agents receive their agent definition file as the
\emph{entire} system prompt, replacing Claude Code's default system
prompt completely \coderef{DOC-1}. This means a sub-agent spawned
via \field{subagent\_type} loses all generic operational guidance
(tool usage patterns, verification discipline, output formatting)
unless that guidance is explicitly provided through the agent
definition or an alternative delivery mechanism.

The current hand-written agent definitions are 40--50 line
personality sketches. They contain no operational baseline, no
extension context (quarry memory, vox config), and no role-specific
enforcement hooks. The result is sub-agents that know \emph{who}
they are but not \emph{how} to work.

Three mechanisms must converge to produce a fully functional
sub-agent:

\begin{enumerate}
  \item \textbf{Agent definition file} ---
    \fpath{.claude/agents/<handle>.md}, loaded at spawn as the system
    prompt. Static, generated by ethos.
  \item \textbf{SubagentStart hook} --- fires at spawn, injects
    \field{additionalContext} with persona, team, and extension
    context. Dynamic, runtime.
  \item \textbf{CLAUDE.md} --- loads separately via normal message
    flow \coderef{DOC-2}. Carries org-specific rules. Not duplicated
    in the agent definition.
\end{enumerate}

% ===================================================================
\section{Claude Code Sub-Agent Architecture}
\label{sec:architecture}
% ===================================================================

This section documents the official Claude Code sub-agent behavior
as of v2.1.80 (March 2026), with references to the official
documentation.

\subsection{What Sub-Agents Receive}

\begin{center}
\begin{tabularx}{\textwidth}{@{}lX@{}}
\toprule
\textbf{Component} & \textbf{Behavior} \\
\midrule
System prompt & Agent definition body replaces the full
  Claude Code default system prompt \coderef{DOC-1} \\
CLAUDE.md & Loads separately via normal message flow; not
  part of the system prompt \coderef{DOC-2} \\
MCP servers & Inherited from parent unless restricted via
  \field{mcpServers} frontmatter \coderef{DOC-3} \\
Skills & \emph{Not} inherited; must be listed in
  \field{skills} frontmatter \coderef{DOC-4} \\
Hooks & \emph{Not} inherited from parent; defined in
  agent frontmatter or via \field{SubagentStart}/\field{SubagentStop}
  hooks in \fpath{settings.json} \coderef{DOC-5} \\
Permissions & Inherited from parent session; overridable via
  \field{permissionMode} except when parent uses
  \field{auto} or \field{bypassPermissions} \coderef{DOC-6} \\
Memory & Optional persistent directory at user, project,
  or local scope \coderef{DOC-7} \\
\bottomrule
\end{tabularx}
\end{center}

\subsection{Hard Constraints}

\begin{enumerate}
  \item Sub-agents \textbf{cannot spawn other sub-agents}
    \coderef{DOC-8}. Nested delegation requires agent teams
    (\field{TeamCreate}).
  \item Plugin-provided agents do \textbf{not} support
    \field{hooks}, \field{mcpServers}, or \field{permissionMode}
    frontmatter \coderef{DOC-9}. Copy to
    \fpath{.claude/agents/} if needed.
  \item Sub-agent transcripts persist independently of
    the main conversation and survive compaction
    \coderef{DOC-10}.
\end{enumerate}

\subsection{Frontmatter Fields}
\label{sec:frontmatter}

\begin{center}
\begin{tabularx}{\textwidth}{@{}llX@{}}
\toprule
\textbf{Field} & \textbf{Required} & \textbf{Description} \\
\midrule
\field{name} & Yes & Unique identifier (lowercase, hyphens) \\
\field{description} & Yes & When Claude should delegate to this
  agent \\
\field{tools} & No & Tool allowlist; inherits all if omitted \\
\field{disallowedTools} & No & Tool denylist \\
\field{model} & No & Model: \texttt{sonnet}, \texttt{opus},
  \texttt{haiku}, full ID, or \texttt{inherit} \\
\field{effort} & No & \texttt{low}, \texttt{medium},
  \texttt{high}, \texttt{max} \\
\field{permissionMode} & No & \texttt{default},
  \texttt{acceptEdits}, \texttt{dontAsk},
  \texttt{bypassPermissions}, \texttt{plan} \\
\field{skills} & No & Skills preloaded at startup (full content
  injected) \\
\field{memory} & No & Persistent memory: \texttt{user},
  \texttt{project}, \texttt{local} \\
\field{hooks} & No & Lifecycle hooks scoped to this agent \\
\field{isolation} & No & \texttt{worktree} for git worktree
  isolation \\
\field{maxTurns} & No & Maximum agentic turns \\
\field{background} & No & Always run as background task \\
\bottomrule
\end{tabularx}
\end{center}

% ===================================================================
\section{Identity Delivery Mechanisms}
\label{sec:delivery}
% ===================================================================

Three mechanisms deliver identity context to a sub-agent. Each
carries different content at different lifecycle points.

\subsection{Mechanism 1: Agent Definition File (Static)}

\begin{description}[style=nextline]
  \item[File] \fpath{.claude/agents/<handle>.md}
  \item[Loaded] At spawn, as the system prompt
  \item[Content] Frontmatter (tools, model, effort, memory, skills,
    hooks) + body (personality, writing style, talents, role
    constraints)
  \item[Generated by] Ethos SessionStart hook
  \item[Regenerated when] Generated content differs from the
    existing file (content comparison)
\end{description}

The agent definition file is \textbf{static} --- it reflects the
state of ethos identity data at the time of generation. It does not
change during a session.

\subsubsection{Body Assembly Order}

The body is assembled from ethos identity attributes in this order
(matching \cmd{buildAgentFile} in
\fpath{internal/hook/generate\_agents.go}):

\begin{enumerate}
  \item \textbf{Opening line} --- first content sentence from
    \fpath{personalities/<slug>.md}, folded into
    ``You are Name (handle), <sentence>.''
  \item \textbf{Reporting line} --- hardcoded
    ``You report to Claude Agento (COO/VP Engineering).''
  \item \textbf{Remaining personality} --- from
    \fpath{personalities/<slug>.md}, after stripping the leading
    heading and first paragraph (already used in the opening line).
  \item \textbf{Writing style} --- from
    \fpath{writing-styles/<slug>.md}, under a \texttt{\#\# Writing
    Style} heading.
  \item \textbf{Responsibilities} --- from the role definition, under
    a \texttt{\#\# Responsibilities} heading. Each responsibility is
    a bullet point.
  \item \textbf{Talents} --- listed inline as comma-separated slugs
    (e.g., ``Talents: engineering, formal-methods'').
\end{enumerate}

\noindent\textbf{Note:} There is no ``What You Do / What You Don't
Do'' section today --- that is future work (Phase~2.1).

\subsubsection{Frontmatter Assembly}

Frontmatter fields are derived from the identity and role:

\begin{center}
\begin{tabularx}{\textwidth}{@{}llX@{}}
\toprule
\textbf{Field} & \textbf{Source} & \textbf{Default} \\
\midrule
\field{name} & \field{identity.handle} & Required \\
\field{description} & Identity + role summary & Required \\
\field{tools} & Role definition & Read, Write, Edit, Bash,
  Grep, Glob \\
\field{model} & Role definition & Omitted when role has no model \\
\field{effort} & (planned, not yet generated) & --- \\
\field{memory} & (planned, not yet generated) & --- \\
\field{skills} & (planned, not yet generated) & --- \\
\field{hooks} & (planned, not yet generated) & --- \\
\bottomrule
\end{tabularx}
\end{center}

\subsection{Mechanism 2: SubagentStart Hook (Dynamic)}
\label{sec:subagentstart}

\begin{description}[style=nextline]
  \item[Hook] \fpath{hooks/subagent-start.sh} $\to$
    \cmd{ethos hook subagent-start}
  \item[Fires] When any sub-agent is spawned
  \item[Output] JSON with \field{additionalContext} injected into the
    sub-agent's context
  \item[Content] Persona block + team context + parent resolution +
    extension \field{session\_context} blocks
\end{description}

The SubagentStart hook provides \textbf{dynamic} context that
reflects runtime state. Unlike the agent definition file, it reads
current extension data, current session roster, and current team
configuration at spawn time.

\subsubsection{Current Behavior}

The existing \cmd{HandleSubagentStart} function
(\fpath{internal/hook/subagent\_start.go}):

\begin{enumerate}
  \item Reads \field{agent\_type} from the hook payload.
  \item If an identity with that handle exists, loads it with full
    attribute resolution.
  \item Calls \cmd{BuildPersonaBlock} to assemble persona context
    (personality, writing style, talents).
  \item Resolves the parent agent from the session roster and adds
    ``You report to Name (handle).''
  \item Calls \cmd{BuildExtensionContext} to collect and emit
    extension \field{session\_context} blocks.
  \item Emits the result as
    \field{hookSpecificOutput.additionalContext}.
\end{enumerate}

\subsubsection{Extension Session Context}
\label{sec:ext-context}

Since v2.6.1 (PR \#152, bead ethos-x4i), the SubagentStart hook
injects extension \field{session\_context} into sub-agents. After
building the persona block, \cmd{HandleSubagentStart}:

\begin{enumerate}
  \item Calls \cmd{BuildExtensionContext(ext)} with the resolved
    identity's extensions.
  \item Iterates over all extension namespaces in
    \fpath{<handle>.ext/*.yaml}.
  \item For each file containing a \field{session\_context} key,
    appends the value to the \field{additionalContext} output.
  \item If the identity has no extensions, skips silently (not an
    error).
\end{enumerate}

This ensures sub-agents receive the same tool-scoped instructions
as the main session --- quarry memory instructions, vox voice
configuration, beadle GPG key configuration, and any future
tool-scoped context --- without baking runtime-mutable extension
data into the static agent definition file.

\subsection{Mechanism 3: CLAUDE.md (Inherited)}

\begin{description}[style=nextline]
  \item[File] \fpath{./CLAUDE.md} or \fpath{./.claude/CLAUDE.md}
  \item[Loaded] Separately, via normal message flow \coderef{DOC-2}
  \item[Content] Org standards, team structure, workflow, tool
    usage, invariants
\end{description}

CLAUDE.md loads into every session --- main and sub-agent ---
automatically. Its content must \textbf{not} be duplicated in agent
definition files or SubagentStart output.

Content that belongs in CLAUDE.md (not in agent definitions):

\begin{itemize}
  \item Team and delegation structure
  \item Biff, beads, quarry, beadle tool sections
  \item Communication rules, banned patterns
  \item Git safety, invariants
  \item Operating principles
\end{itemize}

% ===================================================================
\section{Baseline Operational Skill}
\label{sec:baseline}
% ===================================================================

When the agent definition replaces Claude Code's default system
prompt, the sub-agent loses generic operational guidance. The
\field{skills} frontmatter field provides the official mechanism to
inject content at startup \coderef{DOC-4}: ``The full content of
each skill is injected into the subagent's context, not just made
available for invocation.''

\subsection{Skill: baseline-ops}

\begin{description}[style=nextline]
  \item[Location] \fpath{\~{}/.claude/skills/baseline-ops/SKILL.md}
    (user-level, available in all projects)
  \item[Preloaded via] \field{skills: [baseline-ops]} in every
    generated agent frontmatter
  \item[Content] Operational discipline subset from Claude Code's
    default system prompt ($\sim$50--80 lines)
\end{description}

\subsubsection{Required Content}

\begin{enumerate}
  \item \textbf{Tool usage} --- use Read over \cmd{cat}, Grep over
    \cmd{grep}, Glob over \cmd{find}. Use dedicated tools instead of
    Bash equivalents.
  \item \textbf{Verification} --- run \cmd{make check} after every
    change. Zero violations before returning results.
  \item \textbf{No commits} --- never commit. Return results to the
    COO for commit.
  \item \textbf{Conventions} --- look at neighboring files before
    writing new code. Follow existing patterns.
  \item \textbf{Scope discipline} --- no features beyond the spec.
    No comments unless logic is non-obvious. No docstrings on
    unchanged code.
  \item \textbf{Security} --- never introduce OWASP top~10
    vulnerabilities. Never log secrets.
  \item \textbf{Progress tracking} --- use TodoWrite for multi-step
    work. Mark tasks complete as they finish.
  \item \textbf{Output} --- concise, no preamble, no postamble. Lead
    with the answer.
\end{enumerate}

\subsubsection{What Is NOT in the Baseline}

\begin{itemize}
  \item Org-specific rules (covered by CLAUDE.md)
  \item Personality, writing style, talents (covered by agent
    definition body)
  \item Extension context (covered by SubagentStart hook)
\end{itemize}

% ===================================================================
\section{Persistent Agent Memory}
\label{sec:memory}
% ===================================================================

Claude Code supports per-agent persistent memory via the
\field{memory} frontmatter field \coderef{DOC-7}. When enabled, the
sub-agent receives instructions for reading and writing to its
memory directory, and the first 200 lines of \fpath{MEMORY.md} are
loaded at startup.

\subsection{Scope Selection}

\begin{center}
\begin{tabularx}{\textwidth}{@{}llX@{}}
\toprule
\textbf{Scope} & \textbf{Location} & \textbf{Use Case} \\
\midrule
\texttt{user} & \fpath{\~{}/.claude/agent-memory/<name>/} &
  Cross-project knowledge on one machine \\
\texttt{project} & \fpath{.claude/agent-memory/<name>/} &
  Repo-specific patterns, committable via git \\
\texttt{local} & \fpath{.claude/agent-memory-local/<name>/} &
  Repo-specific, not in git \\
\bottomrule
\end{tabularx}
\end{center}

\subsection{Relationship to Quarry}

Claude Code's \field{memory} and quarry are complementary, not
competing:

\begin{center}
\begin{tabularx}{\textwidth}{@{}lXX@{}}
\toprule
& \textbf{Agent Memory} & \textbf{Quarry} \\
\midrule
Persistence & Git-committable files & Database (local or cloud) \\
Scope & Per-agent, per-repo & Per-agent or global, cross-repo \\
Decay & None --- files persist & Temporal decay for observations,
  permanent for facts \\
Search & Agent reads own files & Semantic search across all
  knowledge \\
Use case & Quick notes, patterns & Deep knowledge, RFCs, design
  docs \\
\bottomrule
\end{tabularx}
\end{center}

Default for all generated agents: \field{memory: project}.

% ===================================================================
\section{Role-Based Hooks}
\label{sec:hooks}
% ===================================================================

Agent frontmatter supports lifecycle hooks scoped to the agent's
execution \coderef{DOC-5}. These make quality gates visible at the
point of the file write instead of leaving them as instructions the
agent may forget to run.

\subsection{Implementation Agents}

Agents whose role includes code modification (rmh, bwk, mdm, adb)
receive a \field{PostToolUse} hook that runs quality gates after
file writes:

\begin{lstlisting}[language={}]
hooks:
  PostToolUse:
    - matcher: "Write|Edit"
      hooks:
        - type: command
          command: "(cd \"$CLAUDE_PROJECT_DIR\" && make check) 2>&1 | head -60"
\end{lstlisting}

This runs \cmd{make check} after every file write and surfaces the
first 60 lines of output in the tool-result stream --- visible
enforcement, not blocking enforcement. The pipe to \cmd{head} masks
the exit code, so Claude Code does not gate the next Write on a
broken build. The hook runs automatically; the agent does not need
to remember to run it, and can still proceed through intermediate
states during a refactor. The window size (60 lines) catches both
the compile-error case --- Go compile errors fire at the top of
\cmd{make check} output and short-circuit the rest --- and the
first-failing-test case, since \cmd{-race -count=1 -v} test runs
surface the first \texttt{FAIL:} near the top before the 500+
lines of trailing verbose output that would push a \cmd{tail -20}
window off the relevant error. The subshell and
\field{\$CLAUDE\_PROJECT\_DIR} expansion pin cwd to the project
root so \cmd{make check} resolves against the repo Makefile even
if the sub-agent has cd'd into a subdirectory before the Write or
Edit tool fires. The command runs under \fpath{/bin/sh} (which is
\cmd{dash} on Debian/Ubuntu); use POSIX-sh syntax only. Bashisms
like \cmd{set -o pipefail}, process substitution, and
\field{\$\{VAR:-default\}} are not portable and must be avoided.

\subsection{Review Agents}

Agents whose role is read-only review (djb for security) do not
receive write hooks. Their tool restrictions
(\field{tools: Read, Grep, Glob, Bash}) already prevent file
modification.

\subsection{Hook Source Priority}

When hooks are defined in both the agent frontmatter and
\fpath{settings.json}, both fire. Agent frontmatter hooks are
scoped to the agent's execution and automatically cleaned up when
the agent finishes \coderef{DOC-5}.

% ===================================================================
\section{Agent Definition Generation}
\label{sec:generation}
% ===================================================================

Ethos generates \fpath{.claude/agents/<handle>.md} for every agent
identity in the current team. Generation runs during the
SessionStart hook.

\subsection{Trigger: SessionStart}

SessionStart is the only reliable trigger. It runs at the start of
every session before any delegation can occur. There is no separate
install step.

\medskip
\noindent\textbf{Note:} Before generation runs,
\cmd{InstallAgentDefinitions} copies any hand-authored \texttt{.md}
files from \fpath{.punt-labs/ethos/agents/} to
\fpath{.claude/agents/}. Generation then writes identity-derived content for team members
to \fpath{.claude/agents/<handle>.md}, replacing any hand-authored
file only when its filename collides with a generated team-member
handle. This two-step sequence ensures both hand-authored and
generated agents are present.

\subsection{Staleness Check}

Before writing, compare generated content against the existing file
(DES-026, content comparison --- the mtime approach was a rejected
alternative):

\begin{enumerate}
  \item Generate the full agent definition content (frontmatter +
    body) in memory.
  \item If \fpath{.claude/agents/<handle>.md} exists, read its
    current content.
  \item If content is identical, skip the write (preserves mtime).
  \item If content differs or the file does not exist, write the new
    content.
\end{enumerate}

\subsection{Generation Scope}

\begin{itemize}
  \item Generate for all team members listed in the repo's
    \fpath{.punt-labs/ethos.yaml} team.
  \item \textbf{Skip} the identity matching the \field{agent} field
    in \fpath{ethos.yaml} --- that is the main session identity,
    handled by SessionStart context injection.
  \item Output to \fpath{.claude/agents/<handle>.md} (project-level).
\end{itemize}

\subsection{Idempotency}

Running generation twice with unchanged sources produces identical
output. No timestamp jitter, no non-deterministic ordering.

\subsection{Full Assembly Example}

For identity \texttt{rmh} with personality \texttt{hettinger},
writing style \texttt{hettinger-prose}, and talent
\texttt{engineering}:

\begin{lstlisting}[language={}]
---
name: rmh
description: "Python specialist on the Punt Labs engineering team"
tools:
  - Read
  - Write
  - Edit
  - Bash
  - Grep
  - Glob
---

You are Raymond H (rmh), Python specialist on the Punt Labs
engineering team.
You report to Claude Agento (COO/VP Engineering).

## Core Principles

There must be a better way. Find it.

- Idiomatic Python over transliterated Java/C
- Use the stdlib --- it exists for a reason
...

## Writing Style

Short sentences, under 30 words. Lead with the answer.
...

## Responsibilities
- Implement Python packages: modules, classes, functions, tests
- Write clean, idiomatic Python following the project's standards
...

Talents: engineering
\end{lstlisting}

\noindent The \field{model} field appears only when the role definition
specifies a model. Fields \field{effort}, \field{memory},
\field{skills}, and \field{hooks} are not yet generated (see
Phase~1 and Phase~4).

% ===================================================================
\section{Team Assessment in the Development Workflow}
\label{sec:team-assessment}
% ===================================================================

Before delegating implementation, the COO assesses whether the
current team has the domain expertise needed. This is an early step
in the dev workflow playbook (see
\fpath{punt-kit/docs/dev-workflow-design.md}).

\subsection{Decision Tree}

\begin{enumerate}
  \item \textbf{Existing specialist covers it} ---
    use as-is. No changes.
  \item \textbf{Existing specialist is close} ---
    extend: add a talent file via \cmd{/ethos:talent create} and
    \cmd{/ethos:talent add <handle> <slug>}. Optionally ingest
    domain references into quarry.
  \item \textbf{No specialist is close} ---
    create a new identity via \cmd{/ethos:identity create} with
    a new personality and talent. Rare.
\end{enumerate}

Option~2 (extend) is the common case. Option~3 (create) applies
only when the domain is genuinely foreign to all current team
members.

\subsection{Quarry Knowledge Seeding}

When extending or creating a specialist for a new domain, ingest
authoritative references into quarry:

\begin{itemize}
  \item Common domains (Python, Go, HTTP) --- model training data is
    sufficient. Add a talent file only.
  \item Niche or internal domains (biff protocol, quarry embedding
    pipeline) --- ingest docs because the model does not know
    internal systems.
  \item Post-training or cutting-edge --- ingest because model
    knowledge may be stale.
\end{itemize}

\subsection{Verification}

After extending or creating an identity, spawn the agent and verify:

\begin{lstlisting}
Agent(subagent_type="<handle>",
      prompt="Who are you? What are your principles?")
\end{lstlisting}

Confirm the identity, personality, talents, and operational baseline
are present.

% ===================================================================
\section{Implementation Phases}
\label{sec:phases}
% ===================================================================

\subsection{Phase 1: Create baseline-ops Skill}
\label{sec:phase1}

\textbf{Status: Designed, not yet implemented.} The operational
guidance described in \S\ref{sec:baseline} currently lives in agent
definition bodies and CLAUDE.md. A standalone
\fpath{\~{}/.claude/skills/baseline-ops/SKILL.md} skill has not yet
been extracted.

\subsection{Phase 2: SubagentStart Extension Context}
\label{sec:phase2}

\textbf{Status: Done --- v2.6.1, PR \#152, bead ethos-x4i.}
\cmd{HandleSubagentStart} now calls \cmd{BuildExtensionContext} to
inject extension \field{session\_context} blocks into sub-agent
context. See \S\ref{sec:ext-context}.

\subsection{Phase 3: Agent Definition Generation}
\label{sec:phase3}

\textbf{Status: Done --- v2.6.0, DES-026.} The SessionStart hook
generates \fpath{.claude/agents/<handle>.md} for every agent
identity in the team. Content comparison staleness check, skip-self
behavior, and idempotency are all implemented in
\fpath{internal/hook/generate\_agents.go}.

\subsection{Phase 4: Role-Based Hook Generation}
\label{sec:phase4}

\textbf{Status: Designed, not yet implemented.} Hook templates per
role category (implementation vs review) are specified in
\S\ref{sec:hooks} but are not yet included in generated agent
frontmatter. Generated agents currently receive no hooks.

\subsection{Phase 5: Rollout}
\label{sec:phase5}

\textbf{Status: Partially complete.} Agent definitions are generated
automatically by the SessionStart hook across repos. The baseline-ops
skill (Phase~1) and role-based hooks (Phase~4) are not yet active,
so generated agents rely on CLAUDE.md and agent definition bodies
for operational guidance and behavioral enforcement.

% ===================================================================
\section{Context Budget Analysis}
\label{sec:budget}
% ===================================================================

Every token in the sub-agent's context costs capacity. The following
estimates guide content sizing.

\begin{center}
\begin{tabularx}{\textwidth}{@{}Xrr@{}}
\toprule
\textbf{Component} & \textbf{Est.\ Lines} & \textbf{Est.\ Tokens} \\
\midrule
Agent definition body (personality + style + talents) & 60--100 &
  400--700 \\
Frontmatter (YAML) & 15--25 & 100--150 \\
baseline-ops skill content & 50--80 & 300--500 \\
SubagentStart additionalContext (persona + extensions) & 30--60 &
  200--400 \\
CLAUDE.md (loaded separately; project-level $\sim$120 lines;
  global CLAUDE.md may be larger) & 100--150 & 500--1000 \\
\midrule
\textbf{Total agent context at spawn} & & \textbf{1500--2750} \\
\bottomrule
\end{tabularx}
\end{center}

This is well within the 200k token context window. The sub-agent
retains $>$95\% of its context for actual work.

% ===================================================================
\section{Documentation References}
\label{sec:refs}
% ===================================================================

\begin{center}
\begin{tabularx}{\textwidth}{@{}lX@{}}
\toprule
\textbf{Ref} & \textbf{Source} \\
\midrule
\coderef{DOC-1} &
  \href{https://code.claude.com/docs/en/sub-agents}{Sub-agents docs}:
  ``Subagents receive only this system prompt (plus basic environment
  details like working directory), not the full Claude Code system
  prompt.'' \\
\coderef{DOC-2} &
  \href{https://code.claude.com/docs/en/sub-agents}{Sub-agents docs}:
  ``CLAUDE.md files and project memory still load through the normal
  message flow.'' \\
\coderef{DOC-3} &
  \href{https://code.claude.com/docs/en/sub-agents}{Sub-agents docs}:
  ``By default, subagents inherit all tools from the main
  conversation, including MCP tools.'' \\
\coderef{DOC-4} &
  \href{https://code.claude.com/docs/en/sub-agents}{Sub-agents docs}:
  ``The full content of each skill is injected into the subagent's
  context, not just made available for invocation.'' \\
\coderef{DOC-5} &
  \href{https://code.claude.com/docs/en/sub-agents}{Sub-agents docs}:
  ``Define hooks directly in the subagent's markdown file. These
  hooks only run while that specific subagent is active and are
  cleaned up when it finishes.'' \\
\coderef{DOC-6} &
  \href{https://code.claude.com/docs/en/sub-agents}{Sub-agents docs}:
  ``If the parent uses auto mode, the subagent inherits auto mode
  and any permissionMode in its frontmatter is ignored.'' \\
\coderef{DOC-7} &
  \href{https://code.claude.com/docs/en/sub-agents}{Sub-agents docs}:
  ``The memory field gives the subagent a persistent directory that
  survives across conversations.'' \\
\coderef{DOC-8} &
  \href{https://code.claude.com/docs/en/sub-agents}{Sub-agents docs}:
  ``Subagents cannot spawn other subagents.'' \\
\coderef{DOC-9} &
  \href{https://code.claude.com/docs/en/sub-agents}{Sub-agents docs}:
  ``Plugin subagents do not support the hooks, mcpServers, or
  permissionMode frontmatter fields.'' \\
\coderef{DOC-10} &
  \href{https://code.claude.com/docs/en/sub-agents}{Sub-agents docs}:
  ``Subagent transcripts persist independently of the main
  conversation.'' \\
\bottomrule
\end{tabularx}
\end{center}

% ===================================================================
\section{Related Beads}
\label{sec:beads}
% ===================================================================

\begin{center}
\begin{tabularx}{\textwidth}{@{}llX@{}}
\toprule
\textbf{Bead} & \textbf{Repo} & \textbf{Description} \\
\midrule
ethos-5po & ethos & Generate agent definitions from identity data
  (completed --- DES-026, v2.6.0) \\
punt-kit-7k7 & punt-kit & Dev workflow playbook with specialist
  delegation \\
punt-kit-kpy & punt-kit & Plugin SessionStart permission audit \\
\bottomrule
\end{tabularx}
\end{center}

\end{document}
